\section*{Розділ VIII. Прикінцеві положення}
\addcontentsline{toc}{section}{Розділ VIII. Прикінцеві положення}

\subsection*{8.1. Порядок внесення змін та доповнень}
\addcontentsline{toc}{subsection}{8.1. Порядок внесення змін та доповнень}
    8.1.1. Право ініціювати внесення змін та доповнень до цього Регламенту мають делегати Конференції студентів Університету (КСУ).

    8.1.2. Пропозиції про внесення змін та доповнень до Регламенту КСУ подаються Спікеру КСУ у письмовій формі з обґрунтуванням необхідності таких змін.

    8.1.3. Перед розглядом КСУ проєкт змін та доповнень до Регламенту підлягає обов'язковій експертизі Студентської уповноваженої делегації (СУД) щодо відповідності Положенню про органи студентського самоврядування та чинним нормативним актам студентського самоврядування.

    8.1.4. Зміни та доповнення до цього Регламенту затверджуються Конференцією студентів Університету (КСУ) простою більшістю голосів від загального складу делегатів КСУ за умови попереднього отримання позитивного висновку СУД.

    8.1.5. У разі надання СУД негативного висновку щодо проєкту змін, вони можуть бути затверджені КСУ лише після їх доопрацювання відповідно до зауважень СУД.

\subsection*{8.2. Тлумачення норм Регламенту}
\addcontentsline{toc}{subsection}{8.2. Тлумачення норм Регламенту}
    8.2.1. Офіційне тлумачення норм цього Регламенту надає Конференція студентів Університету (КСУ) за власною ініціативою або за запитом зацікавлених осіб.

    8.2.2. КСУ має право запитувати у Студентської уповноваженої делегації (СУД) висновки щодо тлумачення окремих норм Регламенту в межах її компетенції, визначеної Положенням про СУД.

    8.2.3. Питання про тлумачення норм Регламенту КСУ розглядаються на засіданнях КСУ у порядку, встановленому цим Регламентом.

    8.2.4. Офіційне тлумачення, надане КСУ, є обов'язковим для застосування всіма учасниками процедур, регламентованих цим документом.

\subsection*{8.3. Набуття чинності}
\addcontentsline{toc}{subsection}{8.3. Набуття чинності}
    8.3.1. Цей Регламент затверджується Конференцією студентів Університету (КСУ) простою більшістю голосів від загального складу делегатів КСУ.

    8.3.2. Регламент Конференції студентів Університету набуває чинності з дня його офіційного оприлюднення на інформаційних ресурсах органів студентського самоврядування Університету після його затвердження КСУ.

    8.3.3. З моменту набуття чинності цим Регламентом втрачають силу всі попередні рішення КСУ щодо процедурних питань роботи КСУ, що суперечать положенням цього Регламенту.

\subsection*{8.4. Співвідношення з іншими нормативними актами}
\addcontentsline{toc}{subsection}{8.4. Співвідношення з іншими нормативними актами}
    8.4.1. Цей Регламент діє у частині, що не суперечить Положенню про органи студентського самоврядування Київського авіаційного інституту.

    8.4.2. У разі виникнення суперечностей між нормами цього Регламенту та Положенням про ОСС, застосовуються норми Положення про ОСС.

    8.4.3. Питання, не врегульовані цим Регламентом, вирішуються відповідно до Положення про органи студентського самоврядування Київського авіаційного інституту та загальних принципів права.

\subsection*{8.5. Перехідні положення}
\addcontentsline{toc}{subsection}{8.5. Перехідні положення}
    8.5.1. Цей Регламент застосовується до поточного складу КСУ з моменту його набуття чинності.

    8.5.2. Засідання КСУ, організація яких розпочалася до набуття чинності цим Регламентом, проводяться за процедурами, що діяли на момент початку їх організації, якщо КСУ не прийме іншого рішення.